\chapter{Many tests, and how to report them}\label{ch:many}
\section{Many tests at once: false discoveries}
\input{generated/nb_hypothesis_tests_9}
\section{Reading any test result}
\input{generated/nb_hypothesis_tests_10}

\section{A checklist for reporting a test}
\begin{enumerate}
\item The question and the null hypothesis, in words.
\item The data: what, where, when, how many observations, and how many of them are independent.
\item The test, and why its assumptions hold, or how they were relaxed.
\item The statistic with its degrees of freedom, for example $t(25.0) = 13.3$.
\item The p-value as a number, not only ``$p<0.05$''.
\item An effect size and a confidence interval, in the units of the data.
\item If many tests were run, how false discoveries were controlled.
\item What the result means for the question, and what it does not show.
\end{enumerate}
